%Begin


%%%%%%%%%%%%%%%%%%%%%%%%%%%%%%%%%%%%%%%%%%%%%%%%%%%%%%%%%%%%%%%
\chapter*{\S \space 34. --- DESCRIPTION HEURISTIQUE PROFINIE DE LA CATÉGORIE DES COURBES ALGÉBRIQUES DÉFINIES SUR DES SOUS-EXTENSIONS FINIES $K$ DE $\overline{\mathbf{Q}}_0/{\mathbf{Q}}$ (I.E. DE ${\mathbf{C}}/{\mathbf{Q}}$))}\thispagestyle{empty}
\addcontentsline{toc}{section}{34. Description heuristique profinie de la catégorie des courbes algébriques définies sur des sous-extensions finies $K$ de $\overline{\mathbf{Q}}_0/{\mathbf{Q}}$ (i.e. de ${\mathbf{C}}/{\mathbf{Q}}$)}
\label{sec:34}
\section*{}

On se borne aux courbes géométriques connexes (par commodité) anabéliennes (par nécessité provisoire, cf. plus bas sur la fa\c{c}on de se débarrasser de cette restriction). La donnée d'un revêtement de $U_{\overline{\mathbf{Q}}_0}$ définit une structure d'extension
$$
1 \to \pi \to \Sigma \to \Gamma^+ \to 1
\leqno(1)
$$
ou encore on a un homomorphisme
$$
E \to \boldsymbol{\Gamma}
\leqno(2)
$$
(image ouverte dans $\boldsymbol{\Gamma}$, de noyau appelé $\pi$) où $\boldsymbol{\Gamma} \subset \boldsymbol{\Gamma}_{{\mathbf{Q}}} \isom \boldsymbol{\Gamma}_{0,3}$ est le sous-groupe ouvert correspondant à $K$. Se donner une telle extension (moyennant $\operatorname{Centre}(\pi) = 1$) revient à se donner une action extérieure de $\boldsymbol{\Gamma}'$ sur $\pi$. Une première question~: faut-il mettre la structure à lacets de $\pi$ dans les données de l'objet (1) (ou (2)) censé décrire $U/K$~?
\vskip .3cm
{
Conjecture. --- \it Ce n'est pas la peine -- la structure à lacets de $\pi$ est la seule structure à lacets invariante par l'action extérieure de $\Gamma$ (ou de $\Gamma^\natural$). Les sous-groupes à lacets $L_i$ sont les sous-groupes maximaux dans $\pi$, tels que $\Norm_E (L_i) \to \Gamma$ ait comme image un sous-groupe ouvert de $\Gamma$\footnote{et tout sous-groupe $\ne \{1\}$ de $\pi$ qui a cette propriété est contenu dans un unique $L_i$\dots}. Je présume aussi que tout homomorphisme entre extensions $E$ de $\Gamma^+$ par un $\pi$, $E'$ de $\Gamma$ par un $\pi'$ ($E$, $E'$ provenant de courbes algébriques $U$, $U'$) et tel que l'image de $E$ dans $E'$ soit ouverte, respecte nécessairement la structure à lacets, et en fait provient d'un homomorphisme (unique, on le sait) de courbes algébriques.
}
\vskip .3cm
En tout cas, si on admet la description des sous-groupes à lacets, il sera clair que l'image par $U$ d'un $L_i$ sera ou bien (1), ou bien contenu dans un unique $L_j$\dots

Ceci signifierait que le foncteur des courbes algébriques sur $K$ (morphismes les morphismes dominants) vers les groupes profinis extérieurs sur lesquels $\Gamma$ opère, serait pleinement fidèle.  Le foncteur ``extension du corps de base'' de $K$ à $K'$ correspond au foncteur restriction d'un groupe extérieur (ou d'une structure d'extension) de $\Gamma$ à $\Gamma'$ (sous-groupe ouvert). Les revêtements étales finis de $U$ correspondent aux $E$-ensembles [$\tilde{U}$ étant choisi]\dots

Pour décrire l'image essentielle de ce foncteur, on n'est pas réduit aux conjectures~: on part de l'extension canonique $E_{0,3}$ de $\boldsymbol{\Gamma}_{0,3}=\boldsymbol{\Gamma}$ par $\hat{\pi_{0,3}}$, on prend un sous-groupe ouvert $E'$ de $E_{0,3}$, d'image $\Gamma \subset \boldsymbol{\Gamma}_{0,3}$, noyau $\pi'\subset \hat{\pi}_{0,3}$, et dans la structure à lacets canonique de $\pi'$ de genre $g$ (induite par celle de $\hat{\pi}_{0,3}$), on prend un $I \subset  I(\pi')$ tel que $2g + \card(I) \ge 3$, stable par $\Gamma$, et on ``bouche les trous'' en $I(\pi') \setminus I$. De même, pour décrire quand une action effective, relevant une action extérieure admissible, est elle-même admissible --- i.e. peut-être obtenue à partir d'une courbe algébrique $U$, \emph{ponctuée} par un point rationnel sur $K$. Ceci ne signifie pas pour autant, que si $U$ est donnée par une action extérieure de $\Gamma$ sur un $\pi$, que les classes de conjugaison de relèvements admissibles de cette action proviennent bien toutes de points de $U$ rationnels sur $K$.  Mais on voit de suite que ceci sera le cas, dès que l'on admet la pleine fidélité pour les \emph{isomorphismes}.

(NB. Même pour les automorphismes de $U_{0,3}=\mathbb{P}^1_{{\mathbf{Q}}}\setminus\{0,1,\infty\}$, cette pleine fidélité n'est pas du tout claire.  Il faudrait prouver que le commutant de $\boldsymbol{\Gamma}_{0,3}$ dans $\hat{\hat{\mathfrak{T}}}_{0,3}$ est réduit à $\mathfrak{S}_3$.  La situation est moins sans espoir, quand on se donne, avec la structure de groupe profini extérieur à opérateur $\Gamma$, une arithmétisation de $\pi$ (ce qui suppose qu'on a explicité une structure à lacets) invariante par $\Gamma$ --- donc dans $\hat{\hat{\mathfrak{T}}}(\pi)$ on dispose d'un $\mathcal{N}(\pi)$, et $\Gamma\hookrightarrow \mathcal{N}(\pi)$. Dans ce point de vue, pour un homomorphisme de $\Gamma$-groupes extérieurs (arithmétisés) $\pi'\to\pi$, il est sous-entendu qu'il est compatible avec les arithmétisations, i.e. si $\pi''$ est l'image de $\pi'$ dans $\pi$, on veut que l'arithmétisation de $\pi''$ déduite de celle de $\pi'$ par ``passage au quotient'', soit celle induite par $\pi$. En particulier, pour les automorphismes extérieurs de $\pi$, il est sous-entendu que non seulement ils commutent à $\Gamma$, mais encore qu'ils sont dans $\mathcal{N}(\pi)$ (ce qui implique qu'ils sont dans $\hat{\mathfrak{T}}(\pi)$, car le centralisateur dans $\pi$ de tout sous-groupe ouvert de $\boldsymbol{\Gamma}_\pi$ est $\{1\}$~!) Dans le cas de $U_{0,3}$, on a $\mathcal{N}_{0,3} \isom \mathfrak{S}_3\times\boldsymbol{\Gamma}_{0,3}$, et on sait que le centre de $\boldsymbol{\Gamma}_{0,3} \isom \boldsymbol{\Gamma}$ est triviale --- donc on trouve bien que le groupe des automorphismes de cette structure est réduit à $\mathfrak{S}_3$~!

Ce point de vue est néanmoins probablement superflu --- car on présume que pour l'action extérieure donnée de $\Gamma$, il y a une unique arithmétisation invariante par $\Gamma$, et que les homomorphismes ``admissibles'' de $\Gamma$-groupes extérieurs ``admissibles'', tels qu'ils ont été définis précédemment, respectent automatiquement cette arithmétisation. S'il n'en était rien, il faudrait bien entendu introduire les arithmétisations dans la structure.

Il se pose la question de trouver une description plus simpliste des actions extérieures d'un $\Gamma$ sur un $\pi$ qui sont ``admissibles''. Ici, on va partir d'un $\pi$ dont on fixe déjà une structure à lacets et une arithmétisation $a$ (invariantes par $\Gamma$), de sorte que $\Gamma\subset \mathcal{N}(\pi)$, $\Gamma\to\boldsymbol{\Gamma}_\pi$ injective à image ouverte, i.e. $\Gamma$ correspond à un scindage partiel (ou germe de scindage) de
$$
1\to \hat{\mathfrak{T}}(\pi)\to \mathcal{N}(\pi)\to
\boldsymbol{\Gamma}_\pi\to 1.
$$
Soit $(g,\nu)$ le type de $\pi$~; si $g\ge 1$ le critère la plus simple, c'est que les ``$\Gamma^\natural$-points'' de $\pi'$ déduits de $\pi$ en bouchant tous les trous, soient distincts. (NB. Si $g=1$, en bouchant tous les trous on tombe dans un cas abélien --- mais \c ca ne fait rien). Si $g=0$, il n'y a pas de condition pour $\nu=3$, et si $\nu>3$, mettant à part une partie $I'\subset  I(\pi)$ avec $\card(I')=3$, la condition c'est qu'en bouchant les trous en $I\setminus I'$, les $\Gamma^\natural$-points de $\pi$ déduits des points de $I\setminus I'$ soient distincts.

On notera que dans cette optique conjecturale, dans les cas limites $(g,0)$ ($g\ge 2$), $(1,1)$, $(0,3)$, on n'impose aucune condition a priori sur les relèvements. C'est peut-être très brutal --- et il se pourrait qu'on trouve des relèvements qui ne correspondent pas à une courbe algébrique --- i.e. qui en fait ne sont \emph{pas} admissibles, même s'ils le paraissent.

Pour y comprendre quelque chose, il me semble qu'il faut revenir à une interprétation des germes de scindages dits ``admissibles'' d'une extension telle que
$$
1\to\hat{\mathfrak{T}}_{g,\nu}\to\mathcal{N}_{g,\nu}\to
\boldsymbol{\Gamma}_{g,\nu}\to 1
$$
où $\boldsymbol{\Gamma}_{g,\nu} \isom \boldsymbol{\Gamma}_{\mathbf{Q}} \isom \boldsymbol{\Gamma}_{0,3}$, comme correspondants aux points algébriques d'une \emph{variété} (plutôt ici, une multiplicité) modulaire $M_{g,\nu,{\mathbf{Q}}}$, dont $\mathcal{N}_{g,\nu}$ est le groupe fondamental arithmétique, et $\hat{\mathfrak{T}}_{g,\nu}$ le groupe fondamental géométrique. J'aimerais examiner cette situation de plus près, par la suite. Pour le moment, il semble prudent de ne pas faire de conjectures hâtives pour une description \emph{directe} (i.e. \emph{pas} via $U_{0,3}$) des opérations extérieures ``admissibles''. On travaillera donc pour le moment avec cette notion sous la forme constructive (via $U_{0,3}$).

Si on a un $\pi$ extérieur de type $(g,\nu)$ avec opération extérieure admissible de $\Gamma\subset \boldsymbol{\Gamma}$, on prévoit qu'il y aura une prédiscrétification stricte canonique sur $\pi$ (pas invariante pas $\Gamma$, bien sûr --- mais telle que l'arithmétisation correspondante le soit). On va même, pour une action effective admissible de $\Gamma^\natural$ sur $\pi$ (relevant l'action extérieure donnée), définir une discrétification orientée correspondante canonique $\pi_0 \subset \pi$ (remplacée par une conjuguée, quand on conjugue le relèvement $\Gamma^\natural\to E$ par un $g\in\pi$) --- toutes ces discrétifications orientées (correspondant aux différents ``points'' de $\B_{\pi,\Gamma^\natural}$) définissent une même prédiscrétification orientée --- et même une même prédiscrétification orientée \emph{stricte}.

L'action effective de $\Gamma^\natural$ sur $\pi$ peut s'obtenir (à isomorphisme près) comme suit~: on prend un sous-groupe ouvert $E'\subset E_{0,3}$, d'où extension $1\to\pi'\to E'\to\Gamma' \to 1$ avec $\pi'\subset \hat{\pi_{0,3}}$, donc $\pi'=\hat{\pi}'_0$ ($\pi'_0=\pi'\cap\pi_{0,3}$).  On a une opération de $\Gamma'$ sur $I(\pi')=I(\pi'_0)$, on la prend triviale (quitte à passer à un sous-groupe ouvert de $\Gamma'$), on choisit $i\in I'\subset  I(\pi')$, on bouche les trous en $I'$, d'où $\pi$, avec discrétification orientée $\pi_0$, et une action extérieure de $\Gamma'$ dessus, qui est relevée en une action effective grâce à $i$.  On trouvera bien ainsi un ``réseau'' --- une discrétification orientée dans $\pi$ --- je dis qu'il ne dépend pas des choix qui ont été faits --- en particulier, que les automorphismes de $(\pi,a,\Gamma^\natural)$ transforment $\pi_0$ en lui-même.  De plus, si on a deux points $x$, $y$ de $\B_{\pi,\Gamma^\natural}$, d'où $\pi(x)$, $\pi(y)$, parmi les classes de chemins de $x$ à $y$ (qui font un bitorseur sous $\pi(y)$, $\pi(x)$) il y a en a pour lesquelles $\pi(x)\to\pi(y)$ envoie $\pi_0(x)$ dans $\pi_0(y)$ --- quand on se limite à ceux-ci, on trouve un sous-groupoïde du groupoïde fondamental $\Pi(\pi,\Gamma^\natural)$, qui est cette fois-ci un groupoïde \emph{connexe} à lacets.  On fera attention que $\Gamma$ n'opère \emph{pas} sur ce groupoïde, bien qu'il opère sur l'ensemble de ses objets.


% End
